\section{DỰ TOÁN VÀ TRIỂN KHAI}

\subsection{Dự toán dự án}

\subsubsection{Căn cứ lập dự toán}
Dự án \gls{geoai} được xây dựng dựa trên các căn cứ pháp lý sau:
\begin{itemize}
    \item \gls{qd-btttt} số 671/QĐ-BTTTT ngày 26/04/2024 về hướng dẫn xác định chi phí phần mềm nội bộ.
    \item Thông tư số 36/2018/TT-BTC ngày 30/3/2018 của Bộ Tài chính về Hướng dẫn việc lập dự toán, quản lý, sử dụng và quyết toán kinh phí dành cho công tác đào tạo, bồi dưỡng cán bộ, công chức, viên chức.
    \item Thông tư số 06/2011/TT-BTTTT ngày 28/2/2011 của Bộ Thông tin và Truyền thông về việc Quy định lập và quản lý chi phí đầu tư ứng dụng \gls{cntt}.
    \item Thông tư 121/2018/TT-BTC ngày 12/12/2018 của Bộ Tài chính về Quy định về lập dự toán, quản lý, sử dụng và quyết toán kinh phí để thực hiện công tác ứng cứu sự cố, bảo đảm an toàn thông tin mạng.
    \item Công văn số 3787/BTTTT-THH ngày 26/12/2014 của Bộ Thông tin và Truyền thông về việc hướng dẫn phương pháp xác định chi phí kiểm thử chất lượng phần mềm.
    \item Công văn số 2519/BTTTT-KHTC ngày 04/9/2014 của Bộ Thông tin và Truyền thông về việc Đơn giá lắp đặt phần cứng và cài đặt phần mềm trong ứng dụng công nghệ thông tin.
    \item \gls{qd-btttt} số 1688/QĐ-BTTTT ngày 11/10/2019 về việc Sửa đổi, bổ sung Quyết định 2378/QĐ-BTTTT về định mức chi phí quản lý dự án, chi phí tư vấn đầu tư ứng dụng \gls{cntt} sử dụng ngân sách nhà nước.
    \item Quyết định số 320/QĐ-BKHCN ngày 30/3/2025 về hướng dẫn xác định đơn giá nhân công.
\end{itemize}

\subsubsection{Dự toán chi tiết}

\textbf{a. Tổng dự toán}

Bảng tổng hợp chi phí phần mềm nội bộ:
\begin{longtable}{|c|p{4.5cm}|p{3.5cm}|c|r|}
\caption{Bảng tổng hợp chi phí phần mềm nội bộ \label{tab:tong_hop_chi_phi}} \\
\hline
\textbf{TT} & \textbf{Khoản mục chi phí} & \textbf{Cách tính} & \textbf{Ký hiệu} & \textbf{Thành tiền (đồng)} \\ \hline
\endfirsthead
\hline
\textbf{TT} & \textbf{Khoản mục chi phí} & \textbf{Cách tính} & \textbf{Ký hiệu} & \textbf{Thành tiền (đồng)} \\ \hline
\endhead
\hline
\multicolumn{5}{r}{{Tiếp theo trang sau}} \\
\endfoot
\hline
\endlastfoot
1 & Chi phí trực tiếp XD, PT, MR phần mềm nội bộ & $G=1,4\times E\times P\times H$ & G & 188.976.132 \\ \hline
2 & Chi phí chung & $G\times 65\%$ & C & 122.834.486 \\ \hline
3 & Thu nhập chịu thuế tính trước & $(G+C)\times 6\%$ & TL & 18.708.637 \\ \hline
4 & Chi phí XD, PT, MR phần mềm nội bộ & $G+C+TL$ & GPM & 330.519.255 \\ \hline
5 & Chi phí dự phòng & Theo quy định & Gdp & 33.051.926 \\ \hline
\multicolumn{3}{|c|}{\textbf{TỔNG CỘNG}} & \textbf{Gpm} & \textbf{363.571.181} \\ \hline
\end{longtable}

Bảng tổng hợp dự toán:
\begin{table}[htbp]
\centering
\caption{Bảng tổng hợp dự toán \label{tab:tong_hop_du_toan}}
\vspace{0.2cm}
\resizebox{\textwidth}{!}{%
\begin{tabular}{|c|p{6cm}|c|p{5cm}|r|r|r|}
\hline
\textbf{Stt} & \textbf{Nội dung chi phí} & \textbf{Ký hiệu} & \textbf{Công thức} & \textbf{Giá trị trước thuế} & \textbf{VAT} & \textbf{Giá trị sau thuế} \\ \hline
I & Chi phí phần mềm nội bộ & GPM & Theo QĐ 671/QĐ-BTTTT & 427.730.802 & 0 & 427.730.802 \\ \hline
II & Chi phí quản lý dự án & GQL & GPM $\times$ tỷ lệ \% & 7.934.406 & 0 & 7.934.406 \\ \hline
III & Chi phí tư vấn đầu tư & GTV & \makecell{GTV1+GTV2 \\ +GTV3+GTV4} & 23.569.401 & 600.000 & 24.169.401 \\ \hline
1 & Chi phí lập báo cáo kinh tế - kỹ thuật & GTV1 & GPM $\times$ tỷ lệ \%, min 10tr & 15.569.401 & 0 & 15.569.401 \\ \hline
2 & Chi phí thẩm tra dự toán & GTV2 & GPM $\times$ tỷ lệ \%, min 2tr & 2.000.000 & 0 & 2.000.000 \\ \hline
3 & Chi phí lập hồ sơ mời thầu & GTV3 & GPM $\times$ tỷ lệ \%, min 3tr & 3.000.000 & 300.000 & 3.300.000 \\ \hline
4 & Chi phí đánh giá hồ sơ dự thầu & GTV4 & GPM $\times$ tỷ lệ \%, min 3tr & 3.000.000 & 300.000 & 3.300.000 \\ \hline
IV & Chi phí khác & GK & & 5.000.000 & 500.000 & 5.500.000 \\ \hline
1 & Chi phí kiểm thử chức năng PM & GKT & Lập dự toán & 0 & 0 & 0 \\ \hline
2 & Chi phí kiểm thử ATANTT & GKTATTT & Lập dự toán & 0 & 0 & 0 \\ \hline
3 & Chi phí thẩm định HSMT & GK1 & GPM $\times$ 0,1\%, min 2tr & 2.000.000 & 200.000 & 2.200.000 \\ \hline
4 & Chi phí thẩm định kết quả đấu thầu & GK2 & GPM $\times$ 0,1\%, min 3tr & 3.000.000 & 300.000 & 3.300.000 \\ \hline
V & Chi phí dự phòng & GDP &  \makecell{(GPM+GQL+GTV+GK) \\ $\times$ 10\%} & 46.423.461 & 0 & 46.423.461 \\ \hline
\multicolumn{2}{|c|}{\textbf{TỔNG CỘNG}} & \textbf{G} & \makecell{GPM+GQL+GTV\\ +GK+GDP} & \textbf{510.658.071} & \textbf{1.100.000} & \textbf{511.758.071} \\ \hline
\end{tabular}
}
\end{table}

\cleardoublepage
\textbf{b. Dự toán chi tiết}

\vspace{0.3cm}
\textbf{- Danh sách các tác nhân (actors):}
\begin{longtable}{|c|p{3cm}|p{6cm}|p{2cm}|c|}
\caption{Danh sách các tác nhân \label{tab:danh_sach_tac_nhan}} \\
\hline
\textbf{TT} & \textbf{Tên tác nhân} & \textbf{Mô tả tác nhân} & \textbf{Phân loại} & \textbf{Trọng số} \\ \hline
\endfirsthead
\hline
\textbf{TT} & \textbf{Tên tác nhân} & \textbf{Mô tả tác nhân} & \textbf{Phân loại} & \textbf{Trọng số} \\ \hline
\endhead
\hline
\multicolumn{5}{r}{{Tiếp theo trang sau}} \\
\endfoot
\hline
\endlastfoot
1 & Cán bộ CNTT (Admin) & Quản lý toàn bộ hệ thống, cấu hình và phân quyền & Phức tạp & 3 \\ \hline
2 & Cán bộ quản lý đô thị & Quản lý tài sản, lập kế hoạch bảo trì, xem báo cáo & Phức tạp & 3 \\ \hline
3 & Công dân & Xem bản đồ công khai, gửi phản ánh sự cố hạ tầng & Phức tạp & 3 \\ \hline
\end{longtable}

\textbf{- Danh sách yêu cầu chức năng:} Bảng ở chương III phần III mục 1. Yêu cầu chức năng của phần mềm.

\vspace{0.3cm}
\textbf{- Bảng chuyển đổi yêu cầu chức năng sang use case:} Bảng ở chương III phần II mục 2. Chuyển đổi yêu cầu chức năng sang Use Case.

\vspace{0.3cm}
\textbf{- Bảng tính điểm tác nhân không hiệu chỉnh (TAW):}
\begin{longtable}{|c|p{4cm}|c|c|c|}
\caption{Bảng tính điểm tác nhân không hiệu chỉnh (TAW) \label{tab:taw}} \\
\hline
\textbf{TT} & \textbf{Loại Actor} & \textbf{Trọng số} & \textbf{Số tác nhân} & \textbf{Điểm} \\ \hline
\endfirsthead
\hline
\textbf{TT} & \textbf{Loại Actor} & \textbf{Trọng số} & \textbf{Số tác nhân} & \textbf{Điểm} \\ \hline
\endhead
\hline
\multicolumn{5}{r}{{Tiếp theo trang sau}} \\
\endfoot
\hline
\endlastfoot
1 & Đơn giản & 1 & 0 & 0 \\ \hline
2 & Trung bình & 2 & 0 & 0 \\ \hline
3 & Phức tạp & 3 & 3 & 9 \\ \hline
\multicolumn{2}{|c|}{Cộng (1+2+3)} & & \textbf{3} & \textbf{9} \\ \hline
\multicolumn{4}{|c|}{\textbf{TAW}} & \textbf{9} \\ \hline
\end{longtable}
\cleardoublepage
\textbf{- Bảng tính điểm các trường hợp sử dụng (use case):}
\begin{longtable}{|c|p{4cm}|c|c|c|c|}
\caption{Bảng tính điểm các trường hợp sử dụng (use case) \label{tab:tbf}} \\
\hline
\textbf{STT} & \textbf{Loại} & \textbf{Trọng số} & \textbf{Hệ số BMT} & \textbf{Số UC} & \textbf{Điểm} \\ \hline
\endfirsthead
\hline
\textbf{STT} & \textbf{Loại} & \textbf{Trọng số} & \textbf{Hệ số BMT} & \textbf{Số UC} & \textbf{Điểm} \\ \hline
\endhead
\hline
\multicolumn{6}{r}{{Tiếp theo trang sau}} \\
\endfoot
\hline
\endlastfoot
1 & B - Đơn giản & 5 & 1 & 11 & 55 \\ \hline
  & B - Trung bình & 10 & 1 & 9 & 90 \\ \hline
  & B - Phức tạp & 15 & 1 & 0 & 0 \\ \hline
2 & M - Đơn giản & 5 & 1.2 & 1 & 6 \\ \hline
  & M - Trung bình & 10 & 1.2 & 0 & 0 \\ \hline
  & M - Phức tạp & 15 & 1.2 & 0 & 0 \\ \hline
3 & T - Đơn giản & 5 & 1.5 & 0 & 0 \\ \hline
  & T - Trung bình & 10 & 1.5 & 1 & 15 \\ \hline
  & T - Phức tạp & 15 & 1.5 & 0 & 0 \\ \hline
\multicolumn{2}{|c|}{Cộng (1+2+3)} & & & \textbf{22} & \textbf{166} \\ \hline
\multicolumn{5}{|c|}{\textbf{TBF}} & \textbf{166} \\ \hline
\end{longtable}
\cleardoublepage
\textbf{- Bảng hệ số phức tạp kỹ thuật – công nghệ (TCF):}
\begin{table}[htbp]
\centering
\caption{Bảng hệ số phức tạp kỹ thuật – công nghệ (TCF) \label{tab:tcf}}
\vspace{0.2cm}
\resizebox{\textwidth}{!}{%
\begin{tabular}{|c|p{5.5cm}|c|c|c|p{3.5cm}|}
\hline
\textbf{TT} & \textbf{Các hệ số KT-CN} & \textbf{Trọng số (Wi)} & \textbf{Giá trị xếp hạng (Fi)} & \textbf{Kết quả ($Wi \times Fi$)} & \textbf{Ghi chú} \\ \hline
\multicolumn{6}{|l|}{\textbf{I. Hệ số KT-CN (TFW)}} \\ \hline
T1 & Xử lý phân tán & 2 & 0 & 0 & 0-5 \\ \hline
T2 & Mức độ quan trọng của hiệu năng & 1 & 2 & 2 & 0-5 \\ \hline
T3 & Hiệu quả sử dụng cho người dùng & 1 & 2 & 2 & 0-5 \\ \hline
T4 & Độ phức tạp của xử lý bên trong & 1 & 2 & 2 & 0-5 \\ \hline
T5 & Khả năng tái sử dụng mã nguồn & 1 & 0 & 0 & 0-5 \\ \hline
T6 & Dễ cài đặt & 0.5 & 1 & 0.5 & 0-5 \\ \hline
T7 & Dễ vận hành & 0.5 & 3 & 1.5 & 0-5 \\ \hline
T8 & Khả năng chuyển đổi & 2 & 0 & 0 & 0-5 \\ \hline
T9 & Dễ dàng bảo trì & 1 & 3 & 3 & 0-5 \\ \hline
T10 & Xử lý đồng thời & 1 & 1 & 1 & 0-5 \\ \hline
T11 & Mức độ hỗ trợ bảo mật & 1 & 2 & 2 & 0-5 \\ \hline
T12 & Sự phụ thuộc vào mã lệnh bên thứ ba & 1 & 1 & 1 & 0-5 \\ \hline
T13 & Mức độ hỗ trợ đào tạo người sử dụng & 1 & 2 & 2 & 0-5 \\ \hline
\multicolumn{4}{|l|}{$TFW = \sum (Wi \times Fi)$} & \textbf{17} & \\ \hline
\multicolumn{6}{|l|}{\textbf{II. Hệ số phức tạp KT-CN (TCF)}} \\ \hline
\multicolumn{4}{|l|}{$TCF = 0.6 + (0.01 \times TFW)$} & \textbf{0.77} & \\ \hline
\multicolumn{4}{|l|}{$UUCP = TAW + TBF$} & 175 & Điểm UC chưa hiệu chỉnh \\ \hline
\multicolumn{4}{|l|}{$AUCP = UUCP \times TCF \times EF$} & \textbf{85.56625} & ĐIỂM USE CASE ĐÃ HIỆU CHỈNH \\ \hline
\end{tabular}%
}
\end{table}
\cleardoublepage

\textbf{- Bảng tính toán hệ số tác động môi trường, nhóm làm việc (ECF):}
\begin{table}[htbp]
\centering
\caption{Bảng tính toán hệ số tác động môi trường, nhóm làm việc (ECF) \label{tab:ecf}}
\vspace{0.2cm}
\resizebox{\textwidth}{!}{%
\begin{tabular}{|c|p{6cm}|c|c|c|}
\hline
\textbf{TT} & \textbf{Các hệ số môi trường} & \textbf{Trọng số (Wi)} & \textbf{Giá trị xếp hạng (Fi)} & \textbf{Kết quả ($Wi \times Fi$)} \\ \hline
E1 & Quen thuộc với quy trình phát triển PM & 1.5 & 5 & 7.5 \\ \hline
E2 & Kinh nghiệm về ứng dụng & 0.5 & 5 & 2.5 \\ \hline
E3 & Kinh nghiệm về hướng đối tượng & 1 & 3 & 3 \\ \hline
E4 & Năng lực chủ trì phân tích & 0.5 & 5 & 2.5 \\ \hline
E5 & Động lực làm việc & 1 & 5 & 5 \\ \hline
E6 & Yêu cầu ổn định & 2 & 5 & 10 \\ \hline
E7 & Nhân viên bán thời gian & -1 & 0 & 0 \\ \hline
E8 & Ngôn ngữ lập trình khó & -1 & 5 & -5 \\ \hline
\multicolumn{4}{|l|}{$EFW = \sum (Wi \times Fi)$} & \textbf{25.5} \\ \hline
\multicolumn{4}{|l|}{$ECF = 1.4 + (-0.03 \times EFW)$} & \textbf{0.635} \\ \hline
\end{tabular}%
}
\end{table}
\cleardoublepage

\textbf{- Bảng tính lương nhân công:}
\begin{table}[htbp]
\centering
\caption{Bảng tính lương nhân công \label{tab:luong_nhan_cong}}
\vspace{0.2cm}
\resizebox{\textwidth}{!}{%
\begin{tabular}{|c|p{4.5cm}|c|p{3.5cm}|r|r|r|}
\hline
\textbf{TT} & \textbf{Hạng mục} & \textbf{Ký hiệu} & \textbf{Cách tính} & \textbf{Kỹ sư bậc 1} & \textbf{Kỹ sư bậc 2} & \textbf{Kỹ sư bậc 3} \\ \hline
1 & Hệ số cấp bậc & HCB & & 2.34 & 2.65 & 2.96 \\ \hline
2 & Hệ số phụ cấp lương & HPC & & 0.00 & 0.00 & 0.00 \\ \hline
3 & Mức lương cơ sở & MLCS & & 2.340.000 & 2.340.000 & 2.340.000 \\ \hline
4 & Lương cơ bản & LCB & HCB $\times$ MLCS & 5.475.600 & 6.201.000 & 6.926.400 \\ \hline
5 & Hệ số điều chỉnh tăng thêm tiền lương & HĐC & & 0.9 & 0.9 & 0.9 \\ \hline
6 & Các khoản đóng góp theo lương & BHLĐ & BHXH + BHYT + BHTN & 1.177.254 & 1.333.215 & 1.489.176 \\ \hline
6.1 & Bảo hiểm xã hội (17,5\%) & BHXH & 17,5\% $\times$ LCB & 958.230 & 1.085.175 & 1.212.120 \\ \hline
6.2 & Bảo hiểm y tế (3\%) & BHYT & 3\% $\times$ LCB & 164.268 & 186.030 & 207.792 \\ \hline
6.3 & Bảo hiểm thất nghiệp (1\%) & BHTN & 1\% $\times$ LCB & 54.756 & 62.010 & 69.264 \\ \hline
7 & Số ngày công trong tháng & t & & 26 & 26 & 26 \\ \hline
8 & Giá ngày công & GNC & \makecell{[(HCB+HPC) $\times$ MLCS \\ $\times$ (1+HĐC) + BHLD] / t} & 445.419 & 504.428 & 563.436 \\ \hline
9 & Giá giờ công & GGC & GNC / 8 & 55.677,38 & 63.053,44 & 70.429,50 \\ \hline
10 & Số nhân công & & & 3 & 0 & 0 \\ \hline
11 & Mức lương lao động bình quân (H) & H & BQ có trọng số theo số NC & 55.677,38 & & \\ \hline
\end{tabular}%
}
\end{table}
\cleardoublepage

\textbf{- Bảng tính toán chi phí trực tiếp xây dựng phần mềm nội bộ:}
\begin{xltabular}{\textwidth}{|c|p{3.5cm}|p{3cm}|c|r|p{3cm}|}
\caption{Bảng tính toán chi phí trực tiếp xây dựng phần mềm nội bộ \label{tab:chi_phi_truc_tiep}} \\
\hline
\textbf{TT} & \textbf{Hạng mục} & \textbf{Diễn giải} & \textbf{Ký hiệu} & \textbf{Giá trị} & \textbf{Ghi chú} \\ \hline
\endfirsthead
\hline
\textbf{TT} & \textbf{Hạng mục} & \textbf{Diễn giải} & \textbf{Ký hiệu} & \textbf{Giá trị} & \textbf{Ghi chú} \\ \hline
\endhead
\hline
\multicolumn{6}{r}{{Tiếp theo trang sau}} \\
\endfoot
\hline
\endlastfoot
\multicolumn{6}{|l|}{\textbf{I. Tính điểm trường hợp sử dụng (Use case)}} \\ \hline
1 & Điểm Actor (TAW) & Phụ lục IV & TAW & 9 & \\ \hline
2 & Điểm Use case (TBF) & Phụ lục V & TBF & 166 & \\ \hline
3 & Tính điểm UUCP & $UUCP = TAW + TBF$ & UUCP & 175 & \\ \hline
4 & Hệ số phức tạp KT-CN (TCF) & $TCF = 0,6 + (0,01 \times TFW)$ & TCF & 0.77 & Phụ lục VI \\ \hline
5 & Hệ số phức tạp môi trường (EF) & $EF = 1,4 + (-0,03 \times EFW)$ & EF & 0.635 & Phụ lục VII \\ \hline
6 & Tính điểm AUCP & $UUCP \times TCF \times EF$ & AUCP & 85.56625 & \\ \hline
\textbf{II} & \textbf{Nội suy thời gian lao động (P)} & Giờ công/điểm UC & P & 20 & Nội suy từ PL VII (20-28 giờ) \\ \hline
\textbf{III} & \textbf{Giá trị nỗ lực thực tế (E)} & $E = 10/6 \times AUCP$ & E & 142.6104167 & Hệ số điều chỉnh nỗ lực \\ \hline
\textbf{IV} & \textbf{Mức lương lao động BQ (H)} & đồng/giờ & H & 55.677,38 & Theo QĐ 320/QĐ-BKHCN \\ \hline
\textbf{V} & \textbf{Chi phí trực tiếp PM nội bộ (G)} & $G = 1,4 \times E \times P \times H$ & G & 222.324.862 & đồng \\ \hline
\end{xltabular}
\cleardoublepage


\textbf{Tổng dự toán}
Bảng tổng hợp chi phí phần mềm nội bộ:
\begin{longtable}{|c|p{4.5cm}|p{3.5cm}|c|r|}
\caption{Bảng tổng hợp chi phí phần mềm nội bộ \label{tab:tong_hop_chi_phi}} \\
\hline
\textbf{TT} & \textbf{Khoản mục chi phí} & \textbf{Cách tính} & \textbf{Ký hiệu} & \textbf{Thành tiền (đồng)} \\ \hline
\endfirsthead
\hline
\textbf{TT} & \textbf{Khoản mục chi phí} & \textbf{Cách tính} & \textbf{Ký hiệu} & \textbf{Thành tiền (đồng)} \\ \hline
\endhead
\hline
\multicolumn{5}{r}{{Tiếp theo trang sau}} \\
\endfoot
\hline
\endlastfoot
1 & Chi phí trực tiếp XD, PT, MR phần mềm nội bộ & $G=1,4\times E\times P\times H$ & G & 188.976.132 \\ \hline
2 & Chi phí chung & $G\times 65\%$ & C & 122.834.486 \\ \hline
3 & Thu nhập chịu thuế tính trước & $(G+C)\times 6\%$ & TL & 18.708.637 \\ \hline
4 & Chi phí XD, PT, MR phần mềm nội bộ & $G+C+TL$ & GPM & 330.519.255 \\ \hline
5 & Chi phí dự phòng & Theo quy định & Gdp & 33.051.926 \\ \hline
\multicolumn{3}{|c|}{\textbf{TỔNG CỘNG}} & \textbf{Gpm} & \textbf{363.571.181} \\ \hline
\end{longtable}

Bảng tổng hợp dự toán:
\begin{table}[htbp]
\centering
\caption{Bảng tổng hợp dự toán \label{tab:tong_hop_du_toan}}
\vspace{0.2cm}
\resizebox{\textwidth}{!}{%
\begin{tabular}{|c|p{6cm}|c|p{5cm}|r|r|r|}
\hline
\textbf{Stt} & \textbf{Nội dung chi phí} & \textbf{Ký hiệu} & \textbf{Công thức} & \textbf{Giá trị trước thuế} & \textbf{VAT} & \textbf{Giá trị sau thuế} \\ \hline
I & Chi phí phần mềm nội bộ & GPM & Theo QĐ 671/QĐ-BTTTT & 427.730.802 & 0 & 427.730.802 \\ \hline
II & Chi phí quản lý dự án & GQL & GPM $\times$ tỷ lệ \% & 7.934.406 & 0 & 7.934.406 \\ \hline
III & Chi phí tư vấn đầu tư & GTV & \makecell{GTV1+GTV2 \\ +GTV3+GTV4} & 23.569.401 & 600.000 & 24.169.401 \\ \hline
1 & Chi phí lập báo cáo kinh tế - kỹ thuật & GTV1 & GPM $\times$ tỷ lệ \%, min 10tr & 15.569.401 & 0 & 15.569.401 \\ \hline
2 & Chi phí thẩm tra dự toán & GTV2 & GPM $\times$ tỷ lệ \%, min 2tr & 2.000.000 & 0 & 2.000.000 \\ \hline
3 & Chi phí lập hồ sơ mời thầu & GTV3 & GPM $\times$ tỷ lệ \%, min 3tr & 3.000.000 & 300.000 & 3.300.000 \\ \hline
4 & Chi phí đánh giá hồ sơ dự thầu & GTV4 & GPM $\times$ tỷ lệ \%, min 3tr & 3.000.000 & 300.000 & 3.300.000 \\ \hline
IV & Chi phí khác & GK & & 5.000.000 & 500.000 & 5.500.000 \\ \hline
1 & Chi phí kiểm thử chức năng PM & GKT & Lập dự toán & 0 & 0 & 0 \\ \hline
2 & Chi phí kiểm thử ATANTT & GKTATTT & Lập dự toán & 0 & 0 & 0 \\ \hline
3 & Chi phí thẩm định HSMT & GK1 & GPM $\times$ 0,1\%, min 2tr & 2.000.000 & 200.000 & 2.200.000 \\ \hline
4 & Chi phí thẩm định kết quả đấu thầu & GK2 & GPM $\times$ 0,1\%, min 3tr & 3.000.000 & 300.000 & 3.300.000 \\ \hline
V & Chi phí dự phòng & GDP &  \makecell{(GPM+GQL+GTV+GK) \\ $\times$ 10\%} & 46.423.461 & 0 & 46.423.461 \\ \hline
\multicolumn{2}{|c|}{\textbf{TỔNG CỘNG}} & \textbf{G} & \makecell{GPM+GQL+GTV\\ +GK+GDP} & \textbf{510.658.071} & \textbf{1.100.000} & \textbf{511.758.071} \\ \hline
\end{tabular}
}
\end{table}

\cleardoublepage
\subsection{Tiến độ triển khai}
\begin{longtable}{|c|p{6cm}|c|c|}
\caption{Tiến độ triển khai \label{tab:tien_do}} \\
\hline
\textbf{Stt} & \textbf{Hạng mục công việc} & \textbf{Thời gian bắt đầu} & \textbf{Thời gian kết thúc} \\ \hline
\endfirsthead
\hline
\textbf{Stt} & \textbf{Hạng mục công việc} & \textbf{Thời gian bắt đầu} & \textbf{Thời gian kết thúc} \\ \hline
\endhead
\hline
\multicolumn{4}{r}{{Tiếp theo trang sau}} \\
\endfoot
\hline
\endlastfoot
1 & Khảo sát yêu cầu, lập hồ sơ thuyết minh thiết kế kỹ thuật, tổng dự toán & 02/2025 & 03/2025 \\ \hline
2 & Phân tích và thiết kế hệ thống & 03/2025 & 03/2025 \\ \hline
3 & Thiết kế giao diện & 03/2025 & 03/2025 \\ \hline
4 & Xây dựng phần mềm & 03/2025 & 06/2025 \\ \hline
5 & Nghiệm thu và đưa phần mềm vào sử dụng & 06/2025 & 06/2025 \\ \hline
\end{longtable}

\subsection{Phương án tổ chức thực hiện}

\subsubsection{Phương án cài đặt, triển khai}
\begin{itemize}
    \item Triển khai trên hạ tầng đám mây lai (hybrid cloud): máy chủ on-premise của TTDL \gls{tptp} cho dữ liệu nhạy cảm, kết hợp cloud công cộng cho khả năng mở rộng.
    \item CI/CD pipeline tự động hóa quá trình build, test và deploy.
    \item Blue-Green Deployment để đảm bảo zero-downtime khi cập nhật phiên bản.
\end{itemize}

\subsubsection{Phương án đào tạo}
Đào tạo cho 3 nhóm đối tượng:
\begin{itemize}
    \item Nhóm 1 - Quản trị viên hệ thống (10 người): Khóa đào tạo 3 ngày tập trung về quản trị hệ thống, cấu hình \gls{gis} server, quản lý dữ liệu không gian, xử lý sự cố. Tài liệu: Hướng dẫn quản trị hệ thống (150 trang).
    \item Nhóm 2 - Cán bộ quản lý đô thị (50 người): Khóa đào tạo 2 ngày về sử dụng \gls{webgis}, tra cứu và quản lý tài sản, xem báo cáo và dashboard. Tài liệu: Hướng dẫn sử dụng (100 trang).
    \item Nhóm 3 - Cán bộ khảo sát thực địa (100 người): Đào tạo 1 ngày về sử dụng mobile app, quy trình thu thập dữ liệu tại thực địa. Tài liệu: Hướng dẫn nhanh (30 trang) + video tutorial.
\end{itemize}
\cleardoublepage
\subsubsection{Phương án vận hành, quản lý, khai thác}
\begin{itemize}
    \item Bố trí 2 cán bộ chuyên trách vận hành hệ thống \gls{geoai} tại Trung tâm \gls{cntt}-TT \gls{tptp}.
    \item Hệ thống giám sát tự động (Prometheus + Grafana) cảnh báo 24/7 về hiệu năng, lỗi và bảo mật.
    \item Họp đánh giá hàng tháng giữa đơn vị vận hành và các sở ngành sử dụng để cải tiến liên tục.
    \item Cập nhật model \gls{ai} định kỳ 6 tháng/lần dựa trên dữ liệu mới thu thập.
\end{itemize}

\subsubsection{Phương án kiểm thử các yêu cầu chức năng và phi chức năng}
\begin{itemize}
    \item Kiểm thử chức năng: Sử dụng công cụ Selenium/Playwright cho automation test \gls{webgis}; Postman cho \gls{api} test; QGIS cho kiểm tra dữ liệu không gian.
    \item Kiểm thử hiệu năng: Apache JMeter mô phỏng 500 người dùng đồng thời, đặc biệt kiểm thử render bản đồ với 100.000+ điểm dữ liệu.
    \item Kiểm thử bảo mật: OWASP ZAP cho quét tự động; chuyên gia bảo mật độc lập thực hiện penetration testing.
    \item Kiểm thử \gls{ai}: Đánh giá độ chính xác model (accuracy, precision, recall, F1-score) trên tập dữ liệu kiểm thử 1.000 ảnh được gán nhãn.
\end{itemize}

\subsubsection{Cam kết bảo hành và hỗ trợ kỹ thuật}
\begin{itemize}
    \item Bảo hành 24 tháng kể từ ngày nghiệm thu, bao gồm sửa lỗi phần mềm, hỗ trợ kỹ thuật và cập nhật bản vá bảo mật.
    \item SLA hỗ trợ: phản hồi trong 2 giờ cho lỗi nghiêm trọng, 8 giờ cho lỗi trung bình, 24 giờ cho yêu cầu thông thường.
    \item Cập nhật model \gls{ai} miễn phí trong 24 tháng bảo hành khi dữ liệu training mới đạt đủ ngưỡng.
    \item Hotline kỹ thuật 24/7 trong suốt thời gian bảo hành.
    \item Sau thời gian bảo hành, đơn vị phát triển cam kết cung cấp dịch vụ bảo trì, nâng cấp theo hợp đồng hỗ trợ riêng.
\end{itemize}
\cleardoublepage
\section*{PHỤ LỤC I - SƠ ĐỒ USE CASE HỆ THỐNG GEOAI}
\addcontentsline{toc}{section}{PHỤ LỤC I - SƠ ĐỒ USE CASE HỆ THỐNG GEOAI}

\subsubsection*{1. Use Case Quản lý tài sản đô thị trên bản đồ}
\begin{itemize}
    \item Tác nhân: Cán bộ quản lý đô thị, Cán bộ khảo sát thực địa, Quản trị viên.
    \item Use case chính: Xem bản đồ, Thêm tài sản mới, Cập nhật tài sản, Tra cứu tài sản, Xuất dữ liệu.
    \item Extend: Upload ảnh tài sản, Phân tích \gls{ai} tự động, Gắn thẻ QR code.
\end{itemize}

\subsubsection*{2. Use Case Module AI và Cảnh báo tự động}
\begin{itemize}
    \item Tác nhân: Hệ thống \gls{ai} tự động, Camera/\gls{iot}, Cán bộ quản lý.
    \item Use case chính: Phân tích ảnh tài sản, Phát hiện sự cố, Gửi cảnh báo, Dự báo bảo trì.
    \item Include: Lưu kết quả phân tích \gls{ai}, Cập nhật tình trạng tài sản, Tạo phiếu công việc.
\end{itemize}

\subsubsection*{3. Use Case Cổng Công dân}
\begin{itemize}
    \item Tác nhân: Công dân/người dân, Cán bộ tiếp nhận phản ánh.
    \item Use case chính: Xem bản đồ công khai, Gửi phản ánh sự cố, Tra cứu tiến độ xử lý.
    \item Extend: Xác thực địa chỉ GPS, Đính kèm ảnh sự cố, Nhận thông báo cập nhật.
\end{itemize}
\cleardoublepage
\section*{PHỤ LỤC II - BẢNG THỐNG KÊ DỮ LIỆU TÀI SẢN ĐÔ THỊ}
\addcontentsline{toc}{section}{PHỤ LỤC II - BẢNG THỐNG KÊ DỮ LIỆU TÀI SẢN ĐÔ THỊ}
Ước tính quy mô dữ liệu tài sản đô thị cần quản lý tại \gls{tptp} (giai đoạn đầu):
\begin{longtable}{|c|p{4cm}|p{3.5cm}|p{5cm}|}
\caption{Bảng thống kê dữ liệu tài sản đô thị \label{tab:phu_luc_2}} \\
\hline
\textbf{STT} & \textbf{Loại tài sản đô thị} & \textbf{Số lượng ước tính} & \textbf{Ghi chú} \\ \hline
\endfirsthead
\hline
\textbf{STT} & \textbf{Loại tài sản đô thị} & \textbf{Số lượng ước tính} & \textbf{Ghi chú} \\ \hline
\endhead
\hline
\multicolumn{4}{r}{{Tiếp theo trang sau}} \\
\endfoot
\hline
\endlastfoot
1 & Cây xanh đô thị & $\sim$450.000 cây & Toàn \gls{tptp} \\ \hline
2 & Đèn chiếu sáng công cộng & $\sim$120.000 bóng & Đường phố và công viên \\ \hline
3 & Biển báo giao thông & $\sim$80.000 biển & Các loại biển báo, hiệu lệnh \\ \hline
4 & Hố ga thoát nước & $\sim$200.000 hố ga & Hệ thống thoát nước \\ \hline
5 & Trụ điện công cộng & $\sim$60.000 trụ & Lưới điện hạ thế \\ \hline
6 & Bảng điện tử quảng cáo & $\sim$5.000 bảng & Màn hình LCD, bảng led \\ \hline
7 & Đường giao thông (phân đoạn) & $\sim$15.000 đoạn & Theo lý trình đường \\ \hline
8 & Cầu, hầm chui & $\sim$1.200 công trình & Cầu các loại \\ \hline
9 & Trạm bơm, cống điều tiết & $\sim$3.000 công trình & Hệ thống thủy lợi nội thị \\ \hline
10 & Camera giám sát công cộng & $\sim$18.000 camera & Đường phố và giao lộ \\ \hline
\end{longtable}
\cleardoublepage
\section*{PHỤ LỤC III - YÊU CẦU HẠ TẦNG TRIỂN KHAI}
\addcontentsline{toc}{section}{PHỤ LỤC III - YÊU CẦU HẠ TẦNG TRIỂN KHAI}
\begin{longtable}{|c|p{4cm}|p{4cm}|p{4.5cm}|}
\caption{Yêu cầu hạ tầng triển khai \label{tab:phu_luc_3}} \\
\hline
\textbf{STT} & \textbf{Thành phần} & \textbf{Cấu hình tối thiểu} & \textbf{Ghi chú} \\ \hline
\endfirsthead
\hline
\textbf{STT} & \textbf{Thành phần} & \textbf{Cấu hình tối thiểu} & \textbf{Ghi chú} \\ \hline
\endhead
\hline
\multicolumn{4}{r}{{Tiếp theo trang sau}} \\
\endfoot
\hline
\endlastfoot
1 & GeoServer (GIS Engine) & 8 CPU, 32GB RAM, SSD 500GB & Có thể cluster 2 node \\ \hline
2 & \gls{ai} Processing Server & GPU NVIDIA A 100/ 16 CPU/64GB RAM & Cho inference \gls{ai} real-time \\ \hline
3 & PostgreSQL + PostGIS & 16 CPU, 64GB RAM, NVMe 2TB & CSDL không gian địa lý \\ \hline
4 & Application Server (API) & 8 CPU, 32GB RAM x 2 (HA) & FastAPI + load balancer \\ \hline
5 & MinIO Object Storage & 4 CPU, 16GB RAM, HDD 10TB & Lưu ảnh tài sản \\ \hline
6 & Redis Cache & 4 CPU, 16GB RAM & Cache bản đồ và phiên \\ \hline
7 & Kafka Message Queue & 8 CPU, 32GB RAM, SSD 500GB & Stream \gls{iot} data \\ \hline
8 & Monitoring Stack & 4 CPU, 8GB RAM & Prometheus + Grafana \\ \hline
\end{longtable}

\vspace{1cm}
\begin{flushright}
- \centering Hết - \\
\centering Đà Nẵng, tháng 6 năm 2025
\end{flushright}